% !TEX root = ../main.tex

\chapter{结论与展望}

\section{本文工作总结}

本文围绕人形机器人遥操作与通用全身运动控制问题，尝试建立从数据采集、动作重定向、控制策略生成到任务级闭环验证的完整技术链路。针对传统方法中数据来源单一、全身控制耦合复杂以及任务接口不统一等问题，本文从系统层面对相关模块进行了整合与设计，并给出了一版面向真实任务场景的论文初稿框架。

在数据采集层，本文构建了由 Noitom 惯性动捕、Nokov 光学动捕与 SONIC 仿真交互数据组成的多源异构动作采集体系。该体系兼顾部署灵活性、采集精度和任务复杂度，为人形机器人获取全身运动数据和移动操作示教样本提供了统一入口。针对人体与机器人之间的形态差异，本文进一步建立了基于 GMR 的动作重定向方法，使人体动作能够映射为 Unitree G1 及全尺寸平台的可执行控制目标。

在控制策略层，本文提出了上下肢解耦的通用全身运动控制框架。下肢部分通过强化学习训练获得前进、后退、原地旋转和左右平移等基础 locomotion 能力，上肢部分通过轨迹接口和逆运动学实现末端目标跟踪，并通过全身协调机制处理重心变化和动作冲突问题。该设计在保留系统模块化优势的同时，兼顾了移动稳定性、操作精度和工程部署可行性。

在应用验证层，本文围绕感知--控制闭环系统开展了两类任务验证。在 CyboRacket 动态球拍任务中，本文重点实现了视觉预测结果到全身控制接口的转换，并验证了机器人在动态目标交互中的全身动作生成能力；在 OpenClaw 语音服务任务中，本文则展示了语义规划、视觉感知和全身控制之间的协同关系，说明所提框架能够服务于多阶段具身任务执行。上述工作共同表明，本文提出的方法具有从遥操作数据采集延伸到自主任务执行的整体潜力。

\section{主要创新点}

\begin{enumerate}
  \item 提出并构建了面向人形机器人全身控制的多源异构人体运动数据采集与统一重定向框架，将惯性动捕、光学动捕和仿真交互数据纳入同一处理链路，提升了动作数据来源的覆盖能力与可用性。

  \item 提出了上下肢解耦的通用全身运动控制架构，将下肢强化学习 locomotion 策略与上肢轨迹跟踪控制有机结合，并通过统一接口实现遥操作与自主任务场景下的控制复用。

  \item 验证了从遥操作数据采集、控制策略生成到感知--控制闭环任务执行的完整技术链路，说明所提方法能够支持动态交互任务与服务型具身任务的系统级集成。
\end{enumerate}

\section{未来工作展望}

尽管本文已经初步构建了较完整的人形机器人全身控制框架，但仍有若干方向值得进一步深入研究。

第一，可进一步探索端到端的全身学习方法。在当前解耦架构下，系统已经具备较好的模块化和可解释性，但上下肢间仍主要依赖接口级信息交互。未来可以结合视觉、语言和动作的大模型框架，将解耦模块置于统一优化视角下进行联合训练，以获得更高层次的动作协同性与任务泛化能力。

第二，可拓展到更复杂的长时序任务场景。当前任务验证主要聚焦动态交互和单轮服务执行，未来可考虑多房间导航、多物体连续操作、长期人机协作等复杂任务。此类任务不仅要求底层控制稳定，还要求高层规划、记忆管理和环境建模具备更强能力，因此需要从系统层进行进一步扩展。

第三，可探索视觉--语言--动作模型与全身控制的深度融合。随着 VLA 模型的发展，人形机器人有望通过统一表示完成从自然语言理解到全身动作生成的跨模态推理。如何将这类高层模型输出稳定地映射到底层可执行控制，并保持真机安全性与实时性，将是未来人形机器人研究的重要方向。

总体而言，本文工作为人形机器人遥操作与通用全身运动控制建立了初步框架，但从实验数据完善、系统稳定性提升到任务泛化扩展，仍有较大研究空间。后续工作将围绕这些方向持续推进，以进一步提升人形机器人在真实世界中的自主作业能力。
